\subsection{Measurement validation}

The current evidence supports three descriptive checks rather than a complete
classifier validation. First, the tracked gold file contains ninety-four
human-reviewed city-platform cases. Ninety of these rows record Codex-assisted
review by the same identified human reviewer, and no case has been coded by an
independent second human. We therefore do not report inter-coder agreement,
category-specific agreement, Cohen's kappa, or adjudication counts. These
quantities require an independent human coder to apply the frozen codebook
without seeing the existing label.

Second, the one-sided surrogate rule produces only nominal-exit labels. The
current screening file contains 203 positive disclosure rows, which collapse
to 158 issuers. Sixty-one issuers also appear in the gold file, and the
surrogate and gold labels are nominal exit for all sixty-one. This quantity is
the observed concordance among post-hoc overlap issuers. Because the overlap
was not selected with known inclusion probabilities, it does not identify
positive predictive value for the full surrogate-positive frame. A Jeffreys
adjustment or Wilson interval changes the treatment of a boundary estimate but
does not correct unknown validation selection.

Third, the current data do not identify recall, calibration, or four-category
error. The source-available candidate pool contains thirty-six named issuers
for which the screen found no direct formal event, but none has an overlapping
gold label. The surrogate output also contains no predicted probabilities and
does not emit substantive-exit, functional-transfer, or liquidation labels.
The current overlap therefore evaluates neither false-negative behavior nor a
four-category confusion matrix.

The next validation stage requires two separate forms of human work. To assess
the one-sided screen, we propose a probability design for the 133 named,
source-available, non-overlap candidate issuers. This frame contains 97
screen-positive issuers and 36 issuers for which the screen found no direct
formal event. Subject to approval of the validation design, we would draw 60
positive issuers by simple random sampling without replacement and review all
36 nonpositive issuers, yielding inclusion probabilities of (60/97) and 1.
Human reviewers would receive the original packets and frozen codebook without
the surrogate output. The resulting design-weighted quantities would identify
positive predictive value and recall within this fixed candidate frame, not
in a national LGFV population. Fifteen source-missing rows with blank issuer
names would remain in a separate coverage audit until their identities and
source packets are resolved.

To assess human-label reliability, a second human coder must independently
code the ninety-four frozen gold packets. The second coder must not see the
existing human or LLM labels before recording the formal event, post-event
function, final label, alternative label, and source references. The two
coders would then adjudicate every disagreement and record which coding rule
produced it. This procedure would permit raw and category-specific agreement,
Cohen's kappa or an appropriate sparse-category alternative, and adjudication
counts within the assembled gold sample. Because the current gold file has no
liquidation case, liquidation-specific agreement would remain unobserved.
